\section{Introduction}

Training and evaluating indoor panorama interpretation systems requires reliable 3D room layout references. A reference specifies both the space to represent and the geometry of its boundaries. Our review of MatterportLayout references for Matterport3D panoramas \cite{chang2017matterport3d,zou2021manhattan} found ambiguities around openings and hidden extensions, together with geometric errors within the selected rooms, including omitted fireplace projections and visible wall turns. Reliable annotation therefore requires an explicit target and careful tracing of the structures that define it.

Even with the target fixed, human annotators produce different layouts. One annotator bridges a fireplace with a straight wall while another traces its recessed sides (Figure~\ref{fig:individual-fireplace}). Annotators who retain the same turn also place its upper and lower endpoints at different positions. These differences raise a connected set of questions: which individual choices recur, how do they enter a collective layout, and how does the result change when we add or replace people? Answering them requires following both structural support and boundary geometry through the aggregation process.

\begin{figure}[t]
\centering
\setcounter{subfigure}{0}
\begin{minipage}{86mm}
\centering
\begin{minipage}[t]{.36\linewidth}
\vspace{0pt}\centering
\ESPanel{\linewidth}{panels/fig01/a.pdf}{Panorama}
\par\vspace{3pt}
\ESPanel{\linewidth}{panels/fig01/b.pdf}{Room context}
\end{minipage}%
\hfill%
\ESPanel{.326\linewidth}{panels/fig01/c.pdf}{A29 (4 corners)}%
\hfill%
\ESPanel{.2840275\linewidth}{panels/fig01/d.pdf}{A59 (8 corners)}

\par\vspace{2pt}
\includegraphics[width=.66\linewidth]{\FigureAssetRoot/panels/fig01/legend.pdf}
\par\vspace{2pt}

\end{minipage}
\caption{(a) Panorama 7y3sRwLe3Va-19 and (b) room context. (c) A29 bridges the fireplace with four corners; (d) A59 traces its sides with eight. The footprints share a scale, with the reference dashed; the lower row shows the room context and the three-dimensional layouts reconstructed from each annotation.}
\label{fig:individual-fireplace}
\end{figure}


Recent work clarifies alternative layout targets and improves reference construction \cite{c2ref06,dong2024panocontext,kanayama2025tof360}, while annotation research models worker differences and aggregates complex responses \cite{c2ref07,c2ref17}. For panoramic layouts, structural choices and positional variation are coupled: omitting a shallow recess changes the wall sequence with little change in floor coverage, whereas shifting a shared upper endpoint changes local wall height without changing the corner count. A study of collective layout quality must measure these differences and preserve their connection to the people who supplied the observations.

We propose a framework for understanding annotation differences and constructing collective panoramic room layouts. It links target-space definition, geometric quality measurement and paired-corner structural consensus: defining which enclosure to annotate, measuring geometric discrepancies, and combining observations into a shared layout. Controlled comparisons vary team size, class composition, specific members and retained teammates, while same-room comparisons examine which individual choices recur across viewpoints. Together, these components connect annotation choices to collective geometry and support reference review and team organization. Our contributions are fourfold:

\textbf{Geometric quality measurement.} We introduce $Q$ to measure spatial extent, boundary placement, orthogonal regularity, and global and local height discrepancies. It detects local height distortions obscured by mean-height prism overlap while remaining invariant to collinear subdivisions that preserve wall geometry. Controlled geometric tests and real annotation defect reports establish the interpretation of these measurements.

\textbf{Paired-corner structural consensus.} We develop a method for constructing a shared layout from multiple annotators' paired upper--lower corner observations. It aligns variable-length cyclic wall sequences, retains turns through distinct-person support, and aggregates their upper and lower positions into a connected layout. This carries shared structural observations and local upper boundaries beyond the floor coverage represented by region voting.

\textbf{Doorway-boundary target definition.} We define doorway-boundary viewpoints, where a camera beneath a doorframe exposes adjoining spaces. Choosing either room or their combined space changes the target enclosure and its closure boundary. Making these alternatives explicit separates target choice from geometric accuracy and provides a basis for annotating and evaluating this distinct scene setting.

\textbf{Group-size and composition effects on collective layouts.} We show that adding annotators affects mean quality and the lower tail differently, while class proportions, specific members and retained teammates shape collective structure and geometry. Controlled within-class and cross-class replacements distinguish the magnitude of change from its direction; teams with identical class proportions can produce different layouts, and retained teammates can reverse the effect of the same replacement. Same-room comparisons further identify recurring structural habits alongside viewpoint-dependent corner choices. These findings inform team composition and targeted review beyond simply increasing the number of annotators.
